\chapter{Agentic Coding}
\label{chapter:AgenticCoding}
% EVOLVE-BLOCK-START

Agentic coding is the practice of directing an AI agent to write, run, and revise code on your behalf.\index{Agentic coding}
It is not passive: the researcher is not a bystander watching the agent write a program.
The researcher is the author who has contracted out the typing.
Every file the agent creates, every function it writes, every import it adds belongs to the researcher and will be submitted, graded, and eventually run again on a clean machine.

This distinction matters because agentic tools can produce large volumes of plausible-looking code very quickly.
The risk is not that the agent writes bad code — it is that the researcher stops reading.
The moment the researcher outsources \emph{comprehension} to the agent, the work ceases to be reproducible, auditable, or defensible.

Section~\ref{section:AgenticCodingLoop} lays out the read--direct--inspect loop; the sections that follow sharpen each phase in turn --- directing the agent with precise instructions, reading the diff it produces (Section~\ref{section:ReadingDiff}), pinning a reproducible environment (Section~\ref{section:ReproducibleEnvironments}), and keeping a clean version history so the work stays auditable.

\section{The Agentic Coding Loop}
\label{section:AgenticCodingLoop}

Agentic coding proceeds in a cycle that alternates between researcher direction and agent execution.\index{Agentic coding loop}

\begin{enumerate}
\item \textbf{Specify.} State clearly what the agent should produce — a function, a script, a test — with explicit constraints on interface, style, and correctness criteria.
\item \textbf{Execute.} The agent writes the code, runs it if appropriate, and reports the result.
\item \textbf{Read.} The researcher reads every change the agent made — every diff, every new file, every modified function.
\item \textbf{Commit or revise.} If the change is correct and understood, commit it. If not, provide corrective direction and return to step~1.
\end{enumerate}

The loop is deliberately incremental.
Each cycle should produce a small, reviewable change — not a thousand-line scaffolding dump.
A change the researcher cannot fully read is a change the researcher cannot own.

\begin{principle}
\label{principle:IncrementalDevelopment}
Prefer small, reviewable changes over large, opaque ones.
An instruction that says ``implement the entire pipeline'' will produce code faster than ``implement the data loader, then stop'' — but the former is harder to verify and harder to correct.
Small steps compound correctly; large steps compound errors.
\end{principle}

\section{Directing the Agent Effectively}

The quality of an agent's code output is strongly predicted by the quality of the researcher's direction.
Vague instructions produce vague implementations; incomplete specifications produce code that solves the wrong problem.

\subsection{Specify the Interface First}

Before asking the agent to implement a function, state its interface: name, parameters, return type, and any invariants.
A well-specified interface constrains the implementation space and makes correctness easier to verify.

\begin{example}
A poor specification:
\begin{lstlisting}[language=python]
# "write a function to score a model"
\end{lstlisting}
A good specification:
\begin{lstlisting}[language=python]
def score(model, X_test, y_test):
    """Return mean squared error of model.predict(X_test) vs y_test.
    Both X_test and y_test are numpy arrays.
    The return value is a non-negative float.
    Raise ValueError if arrays have incompatible shapes."""
\end{lstlisting}
The second version leaves the agent no room to misunderstand the task.
\end{example}

\subsection{Cite the Files That Matter}

When the implementation must interact with existing code, name the files explicitly.
``Extend the loader to support CSV files'' is less precise than ``modify \code{src/loader.py}, specifically the \code{load\_dataset} function, to accept files with a \code{.csv} extension as well as \code{.npz}.''
The agent may otherwise modify a different file, or write new code that duplicates existing functionality.

\subsection{One Task per Instruction}

Compound instructions that ask for several things at once produce compound diffs that are harder to review and harder to revert if one part is wrong.

\begin{remark}
``Add a CSV loader, write tests for it, and update the \code{CLAUDE.md} to document it'' is three tasks.
If the tests are wrong, reverting them also reverts the loader.
Three separate instructions — loader, then tests, then documentation — each produce a single-purpose commit that can be independently reviewed, accepted, or reverted.
\end{remark}

\section{Reading the Diff}
\label{section:ReadingDiff}

A \emph{diff}\index{Diff} is a compact representation of the changes between two versions of a file: lines added are prefixed with~\code{+}; lines removed with~\code{-}.
Reading every diff before committing is the single most important habit in agentic coding.

\begin{example}
A minimal diff showing the addition of a CSV branch to a loader:
\begin{lstlisting}[language=diff]
--- a/src/loader.py
+++ b/src/loader.py
@@ -12,6 +12,10 @@ def load_dataset(path):
     if path.endswith('.npz'):
         data = np.load(path)
         return data['X'], data['y']
+    elif path.endswith('.csv'):
+        df = pd.read_csv(path)
+        X = df.drop(columns=['target']).values
+        y = df['target'].values
+        return X, y
     raise ValueError(f"Unsupported format: {path}")
\end{lstlisting}
Reading this diff reveals that the CSV branch assumes the target column is named \code{'target'} and that the remaining columns are all features.
Neither of these constraints appears in the instruction; both affect correctness on real data.
The diff also introduces a dependency on \code{pandas} (the \code{pd} alias) that must already be present in the environment --- itself a change worth catching in the diff rather than at run time on a clean machine.
\end{example}

Reading diffs develops two habits.
The first is \emph{catching errors}: agents make mistakes — wrong assumptions, missing edge cases, subtle logic errors — that are invisible in the running output but legible in the diff.
The second is \emph{maintaining ownership}: a researcher who reads every diff understands their own codebase; one who does not is dependent on the agent to remember what was done and why.

\begin{remark}
This habit of reading every change and owning the result is what Anthropic's AI-fluency framework names \emph{diligence} --- the discipline of ensuring that the speed an agent offers never costs the researcher responsibility, transparency, or accountability for what remains their own work (Anthropic, 2025).
Its companion competences in that framework --- \emph{description} (the context engineering of Chapter~\ref{chapter:ContextEngineering}) and \emph{discernment} (the evaluation of Chapter~\ref{chapter:Evaluation}) --- are developed elsewhere in these notes under their own names.
\end{remark}

\section{Reproducible Environments}
\label{section:ReproducibleEnvironments}

A piece of code is reproducible when anyone can clone the repository, follow the documented procedure, and obtain the same result.\index{Reproducible environment}
For Python projects, reproducibility requires that the exact versions of all dependencies be recorded and restorable.

\subsection{Dependency Pinning}

A \emph{pinned dependency}\index{Dependency pinning} specifies an exact package version, not a version range.

\begin{example}
A loosely specified dependency and its pinned equivalent:
\begin{lstlisting}[language=toml]
# pyproject.toml — imprecise
[project]
dependencies = ["numpy", "scikit-learn"]

# requirements.txt — pinned
numpy==2.1.3
scikit-learn==1.6.1
scipy==1.14.1
\end{lstlisting}
The loose specification allows a collaborator to install a different version than the author used, potentially changing numerical results.
The pinned specification guarantees the same environment.
\end{example}

\subsection{Virtual Environments}

A \emph{virtual environment}\index{Virtual environment} is an isolated Python installation that contains only the packages a specific project requires.
This prevents conflicts between projects that require different versions of the same package and ensures that the installed packages are exactly the ones listed in the requirements file.

\begin{lstlisting}[language=bash]
python -m venv .venv               # create a virtual environment
source .venv/bin/activate          # activate it (macOS / Linux)
pip install -r requirements.txt    # install pinned dependencies
\end{lstlisting}

Modern projects often use tools like \code{uv} or \code{poetry} to manage virtual environments and pin dependencies automatically.

\subsection{One-Command Run}

The final requirement for reproducibility is a single command that installs dependencies, loads data, runs the pipeline, and reports the result — with no manual steps in between.

\begin{example}
Using \code{just} as a command runner:
\begin{lstlisting}[language=bash]
# justfile
run:
    uv run python -m challenge1

# the researcher (or evaluator) runs:
just run
\end{lstlisting}
If an evaluator cannot reproduce the result with one command, the work is not reproducible.
\end{example}

\section{Clean Version History}

A clean version history is one where each commit has a single purpose, a clear message, and represents a state the project can return to.
An agent that is given broad latitude will often produce long sessions of changes; the researcher's job is to organize those changes into coherent commits.

\subsection{Atomic Commits}

An \emph{atomic commit}\index{Atomic commit} is a commit that contains exactly one conceptual change.
Atomic commits make the history navigable: each commit can be understood, reverted, or cherry-picked independently.

\begin{principle}
\label{principle:AtomicCommits}
Before committing, ask: ``If I need to undo \emph{just this one thing}, can I?''
If the answer is no — because the commit combines two unrelated changes — split it.
A commit history is a journal; each entry should have exactly one topic.
\end{principle}

\subsection{Commit Messages}

A commit message should describe \emph{why} the change was made, not just what changed (the diff already shows what changed).

\begin{example}
Uninformative and informative commit messages for the same change:
\begin{lstlisting}[language=bash]
# poor
git commit -m "fix bug"

# better
git commit -m "Score on test split, not full dataset (was leaking into eval)"
\end{lstlisting}
The second message names the bug and explains why it matters.
Six months later, the author will be grateful for it.
\end{example}

\section*{Further Reading}

\begin{small}
\begin{enumerate}
\item The workshop materials, \code{resources/claude-cli-workshop/} Lessons~01--02 — the primary reference for agentic coding with Claude Code.
\item Chacon, S., and Straub, B., \emph{Pro Git}, 2nd edition, Chapter~2 — git basics and best practices.
\item Kleppmann, M., \emph{Designing Data-Intensive Applications}, Chapter~1 — on reproducibility and reliability as design goals; a formative text for thinking about systems rather than scripts.
\end{enumerate}
\end{small}
% EVOLVE-BLOCK-END
